\clearpage
\subsection{Testing}
Testing checks that the implementation is \emph{correct}, that the code and pipeline do what they are meant to do -- independently of whether the underlying research hypothesis turns out to be true.

\textbf{Data and pipeline tests.} Ingestion will be checked for schema validity, duplicate-free replay after an interrupted collection run, and reconciliation of a sample of collected records against source APIs and blockchain data. A point-in-time check will reject any feature or wallet history that uses information published after its forecast cut-off, since look-ahead leakage is the failure most likely to silently invalidate every downstream result without producing an obvious error. Corporate-action handling (splits, dividends) and timestamp normalization across time zones and daylight-saving changes will also be covered, since these are easy to get subtly wrong and hard to notice later.

\textbf{Calculation tests.} The probability model, wallet-scoring routine and the scenario/risk engine will be unit-tested against reference cases with known outputs: event probabilities of exactly zero and one, a missing asset, a one-sided or illiquid order book, and a singular or near-singular covariance matrix. API-level tests will enforce structural invariants -- portfolio weights non-negative and summing to one, horizons restricted to supported values, VaR and expected shortfall computed consistently with their definitions -- so that a broken invariant fails loudly rather than producing a plausible but incorrect result.

\textbf{Workflow and agent tests.} End-to-end tests will exercise the full path from event selection to a saved, shareable report, including stale or unavailable upstream data. Agent-specific tests will cover malformed or out-of-scope tool requests and adversarial instructions embedded in retrieved market or news text, since retrieved content is treated as untrusted data rather than instructions. The standing requirement is that every number shown to the user, including any number referenced by the AI assistant, traces back to a saved, versioned result rather than to generated text. This will itself be checked by a small set of tests that compare assistant output against the underlying stored result.

\textbf{Indicative metrics.} Pass/fail against the reference cases and invariants above, automated-test coverage of the ingestion and calculation modules, zero duplicate records after a replayed or interrupted collection run, and successful, byte-identical reproduction of a prior run's results from its stored data manifest, configuration and random seed.

\subsection{Evaluation}
Evaluation asks whether the finished system meets the objectives in Section~1.2.

\begin{table}[H]
\smalltable
\caption{How each objective will be evaluated (indicative and subject to refinement after the feasibility audit).}
\begin{tabularx}{\linewidth}{@{}L{20mm} L{42mm} Y@{}}
\toprule\tablehead
\textbf{Objective} & \textbf{Question} & \textbf{Indicative metric} \\
\midrule
Wallet identification & Does the smart-wallet score capture information beyond chance? & Backtested accuracy/Brier score vs.\ a label-shuffled baseline, with sensitivity to the eligibility threshold \\
Wallet-informed probability & Does wallet-weighting improve on the raw market price? & Brier score and log loss across the baseline ladder (raw price $\rightarrow$ +liquidity $\rightarrow$ +anonymous flow $\rightarrow$ +wallet signal) \\
Volatility forecasting & Do event features improve the same regression model? & MSE and QLIKE for the recent-return baseline and its augmented version on identical forecast dates \\
Related-event analysis & Do reviewed related events add information beyond the target event? & MSE and QLIKE for target-only versus target-plus-related-event features on identical forecast dates \\
Agent contribution & Does agent-assisted event selection improve over fixed rules? & Mapping precision and downstream forecast loss for rule-based, agent-assisted and hybrid configurations \\
Dashboard and assistant & Is the system usable, and does it stay evidence-linked? & Formative user study: task completion, numerical agreement between assistant and dashboard, correct abstention on insufficient data \\
\bottomrule
\end{tabularx}
\end{table}

\subsubsection{Experimental design}
Comparisons will use a chronological train/validation/test split, with each event's contracts and snapshots kept together and overlapping return windows purged at split boundaries, so that no test-period information leaks backward. The final test interval will remain untouched until the models and thresholds are fixed. Wallet scores, calibration and scenario inputs will use only information available as of each historical forecast time, matching the point-in-time discipline enforced in testing.

Predictive stability will also be reported across rolling and calendar subperiods. This will reveal whether an apparent improvement is persistent or concentrated in an earlier period. A signal that performs well initially but disappears in later unseen data will be reported as unstable rather than presented as general predictive value.

Given the likely small number of independent announcement dates, results will be reported with confidence intervals rather than as single numbers, and a broader macro sample would require a documented scope revision rather than an informal expansion mid-project. Multiple wallet-level significance tests will use a false-discovery correction rather than being read individually. A null result -- no measurable improvement from wallet or event signals -- is a valid and reportable outcome. Profitability is explicitly not a success criterion.

\subsubsection{Non-functional evaluation}
Beyond whether the numbers are right, the dashboard needs to be evaluated on qualities that determine whether it is actually usable for research in practice.

\textbf{Data freshness.} The interface implies near-real-time market and wallet state, so freshness itself needs to be measured rather than assumed: the lag between a Polymarket price update or a new on-chain fill and its appearance on the Overview and Event detail screens, and the lag between an FOMC announcement and the corresponding event being marked resolved. Data is never truly instantaneous once indexing and reconciliation are involved, so the target is a bounded, disclosed lag rather than zero latency -- the evidence panel already surfaces an as-of time for this reason, and that timestamp's accuracy is itself something to check, not just display.

\textbf{Responsiveness.} Interactive screens (watchlist, event detail, scenario studio) should remain usable while heavier jobs run in the background. Response time will be measured separately for cached/precomputed views versus a fresh scenario run or wallet backfill, since the latter are expected to be slower and are already designed to return a job identifier rather than block the interface.

\textbf{Reliability.} Uptime and error rate of the deployed prototype will be logged informally over the evaluation period (e.g., during the formative study and demonstration), along with the frequency of upstream provider failures (rate limits, stale feeds) and whether the system degrades to a clear "data unavailable" state rather than showing an unlabeled stale number.

\textbf{Cost.} Hosted-model and data-provider cost per research query and per full wallet backfill will be tracked, since this bounds how often scores can realistically be refreshed and how many events the assistant can be exercised against during the evaluation itself.

These are secondary to the correctness and research findings above, and will be reported descriptively (typical and worst-case figures) rather than against hard service-level targets, which would not be meaningful for a research prototype at this scale.

Success requires reproducible data and experiment configuration, a controlled baseline comparison, a tested risk engine, and a usable, evidence-linked interface. Positive, negative and null findings will all be reported. Exact sample sizes, thresholds and the final participant count will be confirmed once the data-feasibility audit (Section~2.1) fixes what history is actually available.
